% p1_construction.tex -- round 403, one lemma:
% P1 as a theorem of the fold construction class.  REFUTED.
% Builds with pdflatex.
% Companion machine check: verify_p1_construction.py.
% Discovery census: experiments/tfpt-discovery/p1_construction_probe.py.
% NO RH CLAIM.  Research documentation, not a theorem of RH.
\documentclass[11pt]{article}

\usepackage[a4paper,margin=2.7cm]{geometry}
\usepackage{amsmath,amssymb,amsthm}
\usepackage{microtype}
\usepackage{xcolor}
\usepackage[colorlinks=true,linkcolor=blue!50!black,urlcolor=blue!50!black]{hyperref}

\newtheorem{lemma}{Lemma}
\newtheorem{prop}{Proposition}
\theoremstyle{remark}
\newtheorem{remark}{Remark}

\newcommand{\indm}{\operatorname{ind}_{-}}
\newcommand{\half}{\tfrac12}
\newcommand{\tr}{\operatorname{tr}}

\title{\bfseries P1 is not a construction theorem\\[2mm]
\large Fixed fold geometry, arbitrary sign-conforming weights:
the class dies}
\author{Stefan Hamann}
\date{August 29, 2026}

\begin{document}
\maketitle

\begin{abstract}\noindent
Round 400 found that dead $\chi$ satisfy
P1${}={}\indm(R_{N-3}-\half I)\le 1$ and that only the
position-scramble breaks it.
The natural suspicion (logged as DCCLXVII) is that P1 is a
\emph{construction} Satz of the fold mask, as the r391
white-block bounds were, so that the whole arithmetic load of
the $R^\dagger$ route would sit in $\mathrm{sch}<0$.
One lemma is tested: P1 holds for every sign-conforming
reweighting of a frozen mask.
It is refuted.
Permute / Gauss / exponential / Dirichlet rescalings of the
von~Mangoldt weights at \emph{fixed} positions give
$\indm$ of order $10$--$25$ on FRAME-A (median-flatten:
$\indm=2$).
A relative $10^{-4}$ jitter already tips a disclosed draw.
The separating invariant is the weight-to-position assignment,
not the sign mask.
No RH claim.
\end{abstract}

\medskip\noindent
\textbf{Status.}\enspace
Lemmas~\ref{lem:interlace}--\ref{lem:scale} are proved
(finite identities).
Proposition~\ref{prop:class} is a finite machine census
that refutes the lemma.
No RH claim.  No $L^*$ claim.  No $R^\dagger\succ\half I$ claim.

\section{The lemma}\label{sec:lem}

Write $R=R_{N-3}$ for the dual Christoffel--Darboux Gram of
the first $N_w-3$ dual OP columns on the hole nodes $Y$, and
$A_0=R-\half I$.
P1 is $\indm A_0\le 1$, equivalently $\lambda_2(R)\ge\half$.
The construction class is: freeze the cosine nodes and the
sign pattern that splits them into $\mu$-atoms and
$\nu$-holes, and vary the positive weights arbitrarily
(permutations within each side, log-Gaussian, exponential,
Dirichlet, median-flatten).
This is the r391 protocol, pointed at $A_0$ rather than at
the Fej\'er energy $B(\Sigma,\Sigma)/D$.
It is \emph{not} the r400 scramble, which redraws
\emph{positions}.

\begin{lemma}[rank-one interlacing]\label{lem:interlace}
A symmetric rank-one update changes $\indm$ by at most one.
Over $\mathbb{Q}$:
$\mathrm{diag}(-1,-1)+e_0 e_0^T=\mathrm{diag}(0,-1)$.
\end{lemma}

\begin{lemma}[PSD need not finish a lift]\label{lem:psd}
A positive-semidefinite update of a negative-definite matrix
may leave every negative eigenvalue negative.
Over $\mathbb{R}$: $\mathrm{diag}(-1,-1)+0.1\,I$ still has
$\indm=2$.
\end{lemma}

\begin{lemma}[scale invariance of $R$]\label{lem:scale}
Joint rescaling $(w_\mu,w_\nu)\mapsto c(w_\mu,w_\nu)$ leaves
$R$ unchanged.
On FRAME-A, $\|R-R(7w)\|_F=3.0\cdot 10^{-13}$.
\end{lemma}

The dual weights of $|cw|$ scale so that the CD Gram on $Y$
is homogeneous of degree zero.
Relative $\mu$/$\nu$ scale is \emph{not} free:
$\mu\mapsto 3\mu$ sends FRAME-A to the vacuous chart
($\indm=0$); $\nu\mapsto 3\nu$ yields $\indm=65$.

\section{The class, refuted}\label{sec:class}

\begin{prop}[construction class]\label{prop:class}
P1 does not hold on the construction class.
On FRAME-A, four seeds of a mask-preserving permutation of
the von~Mangoldt weights give $\indm\in\{20,21,21,22\}$;
log-Gaussian $\{18,19,22,24\}$; exponential $\{11,14,15,11\}$;
Dirichlet $\{11,12,18,12\}$; median-flatten (deterministic)
$\indm=2$.
A disclosed six-draw at relative amplitude $10^{-4}$ already
has one sample with $\indm>1$.
On the sample $\{9,22,33,78,82\}$ MAIN stays at $\indm=1$
and a single permutation yields
$24,30,44,96,102$.
\end{prop}

The sign mask is therefore not the invariant.
Neither is positivity, nor block cardinality, nor joint
scale (Lemma~\ref{lem:scale}).
What separates the P1 worlds from the scramble is the
\emph{assignment} of von~Mangoldt (or $\chi$) masses to the
frozen nodes.
Dead $\chi_3$-$15$ has $\indm=1$ on its true weights and
$\indm=20$ after a permutation of those weights: it sits in
P1 for the same arithmetic reason MAIN does, not because it
shares a construction class with MAIN.

\section{Saturation and the Gram candidate}\label{sec:sat}

The Bernstein--Szeg\H{o} dual on $Y$ (r390) has
$\indm(R^{\mathrm{ref}}-\half I)=49$ at FRAME-A and
$13$ eigenvalues of $R^{\mathrm{ref}}$ at $\half$ to
$10^{-10}$.
Those $13$ modes are not a single exact half-filling that
the source lifts infinitesimally: the count grows
($13,27,38,62,65$ on the sample).
The identity ``$48=49-1$'' is the tautology
$\mathrm{killed}=\mathrm{nneg}_{\mathrm{ref}}-\indm A_0$ on
MAIN, not a class law --- $\mathrm{nneg}_{\mathrm{ref}}$
itself is $49,48,54,124,141$ on the same sample, and a
permutation leaves tens of negatives.

Write $B=R-R^{\mathrm{ref}}$ and
$M=P^{\mathrm{ref}}_- B P^{\mathrm{ref}}_-$.
On MAIN, $M$ is positive definite
($\indm M=0$, $49$ positive eigenvalues, residual $0$,
spectrum in $[0.023,0.710]$): the restriction is a Gram.
Omitting that Gram part returns $\indm=49$.
Keeping \emph{only} that Gram part leaves $\indm=13$
(Lemma~\ref{lem:psd} in situ: the cross terms of the full
$B$ finish the lift down to $1$).
Off the MAIN weights the Gram property dies with P1
(permutation residual $0.14$, $\indm M=6$; scramble
residual $0.15$, $\indm M=8$).
So the Gram observation is a MAIN-weight census, not a
construction Satz and not a proof of P1.

\section{Named kills}\label{sec:kills}

Position-scramble (r400): $\indm=21$.
It breaks P1, but so does a permutation that \emph{keeps}
the mask --- scramble is not the unique kill and is not the
separating invariant.
Two-period $S=21$: $\indm=4$, killed $0$.
Sign-flip of the union weights redraws $Y$ and yields
$\indm\sim 90$.
Omit-Gram mutant: $\indm=49$, as required.
Dead $\chi$ stay at $\indm\le 1$ on their true weights and
leave the class under the same permutation.

\section*{Machine check}

Standalone: rank-one interlacing, PSD-incomplete lift,
threshold $\half$.
Construction: FRAME-A inertia and Gram, scale, $\mu$/$\nu$
split, five-law weight-rand, mild jitter, scramble,
two-period, dead-$\chi$ weights versus permutation.
Discovery census: \texttt{p1\_construction\_probe.py}
(\texttt{--smoke} / full), sample of five rungs plus EXT
$k_z=82$.
Exit line: \texttt{P1 CONSTRUCTION VERIFIED}.

\section*{Claim boundary}

Finite identities for interlacing, an incomplete PSD lift,
and scale invariance of $R$; a named census that P1 fails
as soon as the von~Mangoldt profile is disturbed on a frozen
mask.
No statement about zeros of $\zeta(s)$, in either direction.
No RH claim.

\begin{thebibliography}{9}
\bibitem{r390}
S.~Hamann,
\textit{$G_\varepsilon^\mu$ / Bernstein--Szeg\H{o} Fej\'er},
sealed probe, August 2026.
\bibitem{r391}
S.~Hamann,
\textit{Construction-pure $(R)$ and $(L)$},
standalone note, August 2026,
\texttt{rh/problem/construction\_rl.tex}.
\bibitem{r400}
S.~Hamann,
\textit{Bulk one-defect},
standalone note, August 2026,
\texttt{rh/problem/bulk\_one\_defect.tex}.
\bibitem{r401}
S.~Hamann,
\textit{The augmented edge signature},
standalone note, August 2026,
\texttt{rh/problem/edge\_signature.tex}.
\end{thebibliography}

\end{document}
